\section{Conexión entre cuadripolos}
\subsection{Conexión en cascada}

Para analizar la conexión en cascada entre los cuadripolos 9602 y 9610,
resulta conveniente trabajar con los parámetros de transmisión, ya que la
matriz equivalente se obtiene mediante el producto matricial de las matrices
individuales:

\[
T_{eq}=T_{9602}T_{9610}
\]

A partir de la medición directa de la conexión en cascada, y respetando la
convención adoptada para la matriz de transmisión,

\[
T=
\begin{pmatrix}
A & B\\
C & D
\end{pmatrix},
\]

se obtuvo la siguiente matriz equivalente experimental:

\[
T_{eq,\mathrm{exp}} =
\begin{bmatrix}6,0\angle 81,0^\circ\, & 27,36\angle 87,09^\circ\,\mathrm{mS} \\ 541,18\angle 55,0^\circ\,\Omega & 2,40\angle 61,12^\circ\,\end{bmatrix}
\]

Por otro lado, utilizando las matrices de transmisión calculadas a partir de
las matrices de impedancia de cada cuadripolo, se obtuvo:

\[
T_{eq,calc}=T_{Z,9602}T_{Z,9610}
\]

\[
T_{eq,calc}=
\begin{pmatrix}
34,92\angle83,91^\circ
&
104,80\angle114,74^\circ \ \mathrm{mS}
\\[2mm]
418,29\angle212,26^\circ \ \Omega
&
1,23\angle241,46^\circ
\end{pmatrix}
\]


Al comparar ambas matrices se observan diferencias importantes entre la
medición experimental y el cálculo teórico, en los parámetros \(A\) y \(C\) las fases tienen cierta cercanía, pero los módulos presentan diferencias apreciables. En cambio, en los parámetros
ubicados en la segunda fila aparecen discrepancias mucho más grandes, especialmente en la fase, con diferencias cercanas a \(180^\circ\).

Estas diferencias pueden deberse a errores acumulados durante la medición,
ya que la conexión en cascada se obtiene mediante el producto de matrices
complejas. Por lo tanto, pequeños errores de módulo o fase en las matrices
individuales pueden amplificarse en la matriz equivalente. Además, los
parámetros \(T\) dependen fuertemente de la convención de signos de la
corriente \(I_2\), por lo que una inversión de referencia o una medición de fase
incorrecta puede generar diferencias importantes. También influyen la
resistencia interna de la fuente, las impedancias de los instrumentos y las
conexiones entre cuadripolos.

\subsection{Conexión serie}

Para analizar la conexión serie entre ambos cuadripolos se utilizaron los
parámetros de impedancia \(Z\), ya que en esta configuración la matriz
equivalente se obtiene como la suma de las matrices individuales:

\[
Z_{eq,calc}=Z_{9602}+Z_{9610}
\]

A partir de las matrices medidas individualmente se obtuvo:

\[
Z_{eq,calc}=
\begin{pmatrix}
539,27\angle -21,78^\circ & 250,66\angle -64,77^\circ\\
115,82\angle -69,48^\circ & 378,84\angle -50,57^\circ
\end{pmatrix}
\ \Omega
\]

Por otro lado, la matriz obtenida experimentalmente para la conexión serie fue:

\[
Z_{eq,exp}=
\begin{pmatrix}
546,82\angle -21,80^\circ & 238,72\angle -64,87^\circ\\
113,17\angle -69,51^\circ & 368,16\angle -50,69^\circ
\end{pmatrix}
\ \Omega
\]

Al comparar ambas matrices, se observa una buena concordancia entre los
valores calculados y los valores medidos. Las diferencias en módulo son bajas:
aproximadamente \(1,4\%\) para \(Z_{11}\), \(4,8\%\) para \(Z_{12}\),
\(2,3\%\) para \(Z_{21}\) y \(2,8\%\) para \(Z_{22}\). Además, las diferencias
de fase son muy pequeñas, menores a \(0,2^\circ\) en todos los casos.

Por lo tanto, la conexión serie queda correctamente representada mediante la
suma de las matrices \(Z\) individuales. Las pequeñas discrepancias observadas
pueden atribuirse a errores similares mencionados en el caso anterior.
